% \vspace{-0.3cm}
\section{Conclusions}\label{ch4:sec:Conclusion}
This paper addresses the critical challenge of power consumption and performance management within datacenter and HPC cluster nodes by presenting a dynamic power adjustment approach that seeks to minimize energy use while preserving application performance. Leveraging offline reinforcement learning, our strategy uses datasets collected from static characterization of compute nodes across various benchmarks to train an effective RL agent that dynamically optimizes hardware configuration ($PCAP$).

Experimental validations demonstrate our system's robustness, yielding standard deviations in execution time ranging from a minimum of 0.47\unit{s} to a maximum of 2.91\unit{s} and for energy consumption from a minimum of 0.01\unit{kJ} to a maximum of 1.22 \unit{kJ}, indicating consistent energy savings across both memory-bound and compute-bound benchmark tasks.
Remarkably, the controller proved its efficacy when applied to new and phase-varied applications, achieving an average energy saving of 20.34\% with only a minimal average performance degradation of 7.4\%. Repeated tests have further solidified the controller's capability to significantly enhance energy efficiency. These findings underscore the potential of RL-based control for overcoming the limitations of application- and hardware-centric power management designs. Our research bears particular importance for cloud service providers and HPC cluster administrators aiming to significantly reduce energy expenditure with only a marginal compromise on performance. In the future, we intend to explore the deployment of our algorithms in GPU-intensive tasks and wider cluster job environments. 
%The promising results open new avenues for sustainable computing practices, potentially revolutionizing energy management in datacenters.

While our approach has successfully reduced energy consumption with minimal impact on performance, we have so far limited our use of progress as an application-agnostic signal for online performance to applications that can exhibit an iterative behavior, with a heartbeat that occurs often enough during a control interval and is somewhat stable in rate for a given kernel. Future work will concentrate on widening this approach, from extracting a stable signal for a wider set of metrics (like OpenMP callbacks or MPI activity) and dealing with applications than iterate over a heterogeneous workload. 
%This next stage of research will continue to emphasize balancing execution time constraints with the overarching goal of energy-efficient computing, ensuring that optimal performance is achieved within the required timeframe without compromising on energy savings.

% Akhilesh, add github repo
The code and experimental data are available in open source in our \href{https://github.com/akhileshraj91/summer2024.git}{GitHub} page.

% This paper addresses the challenge of power consumption and management in data center HPC nodes by proposing a dynamic power adjustment approach that minimizes energy consumption while maintaining performance. Our approach leverages offline reinforcement learning, using pre-collected datasets generated during the static characterization of an HPC node running different benchmark applications. The trained RL-agent effectively controls the hardware to optimize performance. Experimental validation showed results for systems running memory-bound and compute bound benchmark applications. The controller also displayed improved performance when deployed on new applications and application with varying phases. The repeatability tests on different HPC nodes showed improved energy efficiency with negligible values of performance. These findings suggest that RL-based approaches can solve the problem of application and hardware centric design of power controllers. Our proposed approach is particularly relevant for cloud providers and high-performance cluster administrators aiming to reduce energy consumption with negligible compromise on application performance. Future work will focus on testing the proposed algorithms on GPU workloads and cluster jobs.


% \section*{Acknowledgments}
% \todo{ack GPT if used to help, and specify which parts of the text}
% \todo{ack ChameleonCloud}